\documentclass[11pt]{article}
\usepackage[a4paper,margin=1in]{geometry}
\usepackage{amsmath,amssymb,amsthm,booktabs}
\usepackage[T1]{fontenc}
\usepackage{lmodern}
\usepackage[hidelinks]{hyperref}
\setlength{\emergencystretch}{3em}

\theoremstyle{plain}
\newtheorem{theorem}{Theorem}
\newtheorem{lemma}[theorem]{Lemma}
\newtheorem{proposition}[theorem]{Proposition}
\theoremstyle{definition}
\newtheorem{definition}[theorem]{Definition}
\theoremstyle{remark}
\newtheorem{remark}[theorem]{Remark}

\title{A Proof of Conjecture 00000000035:\\
A Monochromatic Schur Triple with $xy+1$ Prime}
\author{%
  Yinuo Liu\thanks{Independent researcher (\texttt{champagnepapihz}).}
  \and
  opencode-agent\thanks{Autonomous agent operated by the first author.}
  \and
  Muse Spark\thanks{Meta.}%
}
\date{2026-10-04}

\begin{document}
\maketitle

\begin{abstract}
Conjecture 00000000035 asserts that there exists $N$ such that any 2-coloring
of $[N]=\{1,\dots,N\}$ contains a monochromatic Schur triple $x+y=z$ with
$xy+1$ prime. We prove this with the explicit witness $N=7$: with only
$2^7=128$ colorings and 12 Schur triples (allowing $x=y$), the claim is a
finite check, confirmed by independent Python brute force and formalized in
Lean~4 against Mathlib by \texttt{decide}.
\end{abstract}

\section{Statement and reading}

Conjecture 00000000035 states (English original):
\begin{quote}
There exists $N(k)$ such that any 2-coloring of $[N]$ contains a monochromatic
Schur triple $(x,y,z)$ with $x+y=z$ and $xy+1$ prime.
\end{quote}

We adopt the standard readings, recorded here so the formalization matches:
\begin{itemize}
  \item $[N]=\{1,\dots,N\}$ with the usual 1-based indexing.
  \item A 2-coloring is a map $c:[N]\to\{0,1\}$.
  \item A Schur triple is $x,y,z\in[N]$ with $x+y=z$; $x=y$ is allowed
    (e.g.\ $2+2=4$), per the standard combinatorics convention. The conjecture
    does not require $x\ne y$.
  \item Monochromatic means $c(x)=c(y)=c(z)$.
  \item Primality is over $\mathbb{N}$ in the usual sense ($xy+1$ prime).
  \item The symbol $N(k)$ binds no $k$ anywhere else in the statement; $k$ is
    vestigial. We read ``there exists $N(k)$'' as ``there exists $N$''.
\end{itemize}

\begin{theorem}\label{thm:main}
Let $N=7$. Every 2-coloring of $[7]$ contains a monochromatic Schur triple
$x+y=z$ with $xy+1$ prime.
\end{theorem}

\section{Proof}

There are $2^7=128$ colorings (64 up to global color swap, fixing $c(1)=0$),
and 12 Schur triples with $x\le y$ in $[7]$:
\begin{equation}\label{eq:triples}
\begin{aligned}
&(1,1,2),\ (1,2,3),\ (1,3,4),\ (1,4,5),\ (1,5,6),\ (1,6,7),\\
&(2,2,4),\ (2,3,5),\ (2,4,6),\ (2,5,7),\ (3,3,6),\ (3,4,7).
\end{aligned}
\end{equation}
Of these, the values of $xy+1$ are
$2,3,4,5,6,7,5,7,9,11,10,13$ respectively, so all except
$(1,3,4)$, $(1,5,6)$, $(2,4,6)$, $(3,3,6)$ are \emph{good} (i.e.\ $xy+1$
prime). The claim is that every coloring makes at least one of the 8 good
triples monochromatic.

This is verified by exhaustive enumeration: for each of the 64 colorings with
$c(1)=0$, at least one good triple in~\eqref{eq:triples} is monochromatic.
The check was run independently in Python (sieve primes, enumerate all Schur
triples with $x\le y$, test all $2^{N-1}$ colorings with $c(1)=0$):
\begin{itemize}
  \item $N=5$: 1 avoider up to global color swap (2 in total),
    e.g.\ $c=(0,1,0,0,1)$;
  \item $N=6$: 2 avoiders up to global color swap (4 in total),
    e.g.\ $c=(0,1,0,0,1,0)$;
  \item $N=7$: 0 avoiders --- every coloring contains a good monochromatic
    triple.
\end{itemize}
Since an avoider on $[N]$ restricts to an avoider on every smaller interval,
$N=7$ having 0 avoiders implies 0 avoiders for all $N>7$.
Hence $N=7$ is the least witness, proving Theorem~\ref{thm:main} and therefore
the conjecture. \qed

\begin{remark}[Sharpness]
$N=7$ is best possible: $[5]$ and $[6]$ admit colorings with no good
monochromatic triple (see the avoiders above), and the $N=6$ avoider restricts
to an avoider on every smaller $N$, so no $N\le6$ works.
\end{remark}

\section{Formal verification}

The accompanying Lean~4 project formalizes Theorem~\ref{thm:main} against
Mathlib, using Mathlib's \texttt{Nat.Prime} throughout. Colors are
\texttt{Fin~2}; positions are \texttt{Fin~7} with value \texttt{i.val+1}
(the $+1$ bridges the 0-based \texttt{Fin} index and the 1-based $[7]$).

\begin{itemize}
  \item \texttt{schur\_prime\_7} --- Theorem~\ref{thm:main}: for every
    \texttt{c : Fin 7 $\to$ Fin 2} there exist \texttt{x y z : Fin 7} with
    \texttt{(x.val+1)+(y.val+1) = z.val+1}, \texttt{c x = c y},
    \texttt{c y = c z}, and
    \texttt{Nat.Prime ((x.val+1)*(y.val+1)+1)}, discharged by
    \texttt{decide} (with \texttt{maxRecDepth 100000} and
    \texttt{maxHeartbeats 2000000} since the default recursion limit is hit
    when enumerating $2^7$ colorings).
  \item \texttt{conjecture\_00000000035} --- the existential
    \texttt{$\exists$ N, $\forall$ c : Fin N $\to$ Fin 2, $\dots$} witnessed by
    \texttt{$\langle$7, schur\_prime\_7$\rangle$}.
\end{itemize}

The project builds with \texttt{lake build} on
\texttt{leanprover/lean4:v4.33.1} with Mathlib \texttt{v4.33.1}. It contains
no \texttt{sorry}, no \texttt{admit}, no \texttt{native\_decide}, and
introduces no axioms of its own:
\texttt{\#print axioms conjecture\_00000000035} reports exactly the three
standard axioms \texttt{propext}, \texttt{Classical.choice} and
\texttt{Quot.sound}.

\end{document}
